% ── Chapter 1: Introduction ──────────────────────────────────

\chapter{Introduction}

\section{Background and Recent Research}

Environmental monitoring has always mattered, but the stakes have grown
considerably in recent years. Rising concerns around indoor air quality,
urban heat, and industrial emissions have pushed the need for
\emph{continuous, real-time} data collection well beyond what
traditional fixed-station systems can offer. Parameters such as
temperature, humidity, and CO$_2$ concentration are not just numbers on
a screen; they inform public health decisions, drive building
automation, and shape environmental compliance. The trouble is that
most deployments today anchor a sensor to a wall and call it done,
which works reasonably well for a rooftop weather station but fails
the moment conditions vary across space, as they almost always do in
real buildings and work environments.

IoT has changed what is practically achievable.
Cheap microcontrollers, affordable wireless modules, and accessible
cloud platforms mean that a researcher or engineer can now build a
system that reads sensor data every few seconds, stores it locally
in a database, and pushes it to a cloud backend accessible from
anywhere in the world, all for a fraction of what this would have
cost a decade ago \cite{r1,r2}. IoT-based pollution monitors, indoor
air quality dashboards, and distributed sensor networks have all
benefited. The persistent limitation, however, is that these systems
remain static. A Raspberry Pi sitting on a shelf does excellent work
as a fixed monitor, but it cannot tell you what the air quality
looks like five metres to the left.

Communication infrastructure is the glue that holds distributed
monitoring together. Systems that cannot reliably transmit their data
are not particularly useful, however good their sensors are. A lot of
research has gone into vehicular communication protocols, secure
data exchange, and network architecture for distributed systems
\cite{r3,r4}, and those foundations carry over directly to IoT
monitoring applications. HTTP-based communication through a lightweight
framework like Flask has proven to be a practical middle ground for
embedded deployments, being simpler to set up than a full MQTT broker,
and sufficient for the polling intervals that environmental monitoring
typically requires.

In mobile robotics, the story has followed a somewhat separate track.
Foundational work in vision-guided navigation showed that image-based
control was feasible even with the hardware of the late 1990s
\cite{r5}. Subsequent research on homography-based visual servoing
and adaptive trajectory tracking extended these ideas into richer
control scenarios with experimental validation \cite{r6,r7,r8}.
Control strategies for differential-drive platforms (PID, Model
Predictive Control, and fuzzy logic) have all been tested and compared,
with PID consistently emerging as the preferred choice for structured
indoor environments because it is straightforward to tune and
well-understood \cite{r9,r10,r11}.

The perception side of robotics has been transformed by deep learning.
Instance segmentation models can now detect and outline paths, lanes,
and obstacles in real time, even on embedded processors without a
dedicated GPU \cite{r12,r13}. The challenge has shifted from
\emph{whether} vision-based navigation is possible to \emph{how
efficiently} it can run on low-power hardware.

The gap that persists across all of this is integration. Most research
picks one thread (sensing, communication, or navigation) and develops
it carefully while treating the others as secondary. A system that
simultaneously moves through a space, collects environmental data, and
makes that data available in real time across multiple platforms has
received comparatively little attention. This work is an attempt to
close that gap in a practical and reproducible way.

\section{Motivation}

A single CO$_2$ sensor placed in one corner of a room tells you
nothing about what conditions are like near the other corner, especially
in large or open-plan spaces where ventilation is uneven. Achieving
real spatial coverage with fixed sensors requires a dense and expensive
grid of monitoring points. Mounting the sensor on a mobile robot that
drives a repeatable path through the space is a far simpler and more
cost-effective alternative, since the coverage goes where the robot goes.

The motivation here extends beyond mobility alone. The design
philosophy throughout this project was to keep the two subsystems,
namely the robot and the environmental monitor, genuinely independent of each
other. The monitoring system operates correctly whether or not it is
sitting on a robot; the robot operates correctly whether or not it is
carrying sensor hardware. This separation is intentional. It makes
each part easier to develop, test, and maintain in isolation, and it
means the combined system can be extended in either direction without
breaking the other half.

Flask-based HTTP communication was chosen as the interface layer for
the same practical reasons. It is familiar, easy to debug, and
integrates naturally with Firebase, ThingSpeak, and any browser-based
dashboard. There is no dependency on proprietary APIs or dedicated
message brokers. Anyone with a Raspberry Pi and a network connection
can replicate it.

The simplest statement of the motivation is this: to build a complete
and working system that combines a robot that follows a path
autonomously, a sensor unit that reads temperature, humidity, and air
quality in real time, and a communication layer that makes all of that
data accessible both locally and on the cloud, and to validate every
part of it experimentally.

\section{Problem Statement}

Fixed environmental monitoring stations are limited to a single
location. For many practical scenarios such as warehouse air quality surveys,
greenhouse inspections, and multi-zone occupancy monitoring, a static
reading captures only a small slice of the full picture. Environmental
conditions vary spatially, sometimes significantly, and a single fixed
sensor can produce data that is accurate for its exact location but
unrepresentative of the space as a whole.

Mobile robotic platforms offer a way around this, but most existing
approaches either couple sensing and navigation so tightly that the
system is hard to modify, or treat them as fully separate concerns that
are never actually combined into a working platform. Research on
communication focuses on network-level performance metrics \cite{r4},
while robotics research concentrates on navigation and control
algorithms \cite{r8,r10}. The two areas rarely meet in a single,
deployed, experimentally validated system.

Data accessibility is another unresolved issue. A robot collecting
environmental data that is available only through a direct USB
connection or stored on a local SD card misses most of the practical
value. The data needs to be queryable from any device, in real time,
without requiring specialist software on the client side.

These gaps define the central problem this work addresses: the absence
of a modular, integrated, and practically deployable framework that
combines autonomous mobile navigation, real-time environmental sensing,
and multi-platform cloud communication in a single working system.

\section{Objectives of the Study}

The primary objectives of this research are:

\begin{enumerate}
    \item To design and develop a real-time environmental monitoring
          system capable of acquiring temperature, humidity, and air
          quality data from sensor hardware.
    \item To integrate the monitoring system with a mobile robotic
          platform, enabling dynamic data collection across a physical
          space.
    \item To implement HTTP-based communication using a Flask server
          for real-time data transmission and remote accessibility.
    \item To enable simultaneous multi-platform data storage by
          integrating:
          \begin{itemize}
              \item SQLite for on-device local storage
              \item Firebase Realtime Database for cloud-based remote
                    accessibility
              \item ThingSpeak for IoT data visualisation and charting
          \end{itemize}
    \item To validate the complete framework experimentally, covering
          both autonomous navigation performance and continuous
          environmental data acquisition.
\end{enumerate}

\section{Scope of the Study}

This study covers the design and prototype-level implementation of the
full framework. On the sensing side, this includes reading temperature,
humidity, and air quality from an Arduino-interfaced sensor pair,
processing and classifying those readings on a Raspberry Pi~3, and
distributing the data across SQLite, Firebase, ThingSpeak, and a
live Flask dashboard simultaneously.

On the navigation side, the scope covers training a custom YOLOv8n-seg
instance segmentation model on a red-tape indoor path dataset,
optimising it for embedded deployment via ONNX conversion, and running
a PID-based steering controller in a four-thread concurrent software
architecture on a Raspberry Pi~5. Experimental evaluation uses 20
manual teleoperation trials and 20 autonomous navigation trials on a
fixed rectangular indoor test path.

The study does not extend to outdoor environments, curved dynamic paths,
multi-robot coordination, or sensor types beyond the DHT11 and MQ135.
The communication mechanism is HTTP over a local Wi-Fi network;
cellular connectivity is outside the scope. The result is a validated
prototype that demonstrates the feasibility of the proposed approach,
rather than a production-hardened system.

\section{Organisation of the Thesis}

The remainder of this thesis is organised as follows:

\begin{itemize}
    \item \textbf{Chapter~2} reviews existing literature covering
          mobile robotics, environmental monitoring, control strategies,
          and communication frameworks. It identifies the research gaps
          that motivated this work.
    \item \textbf{Chapter~3} describes the complete system methodology:
          hardware assembly, dataset collection, model training and
          optimisation, autonomous navigation implementation, and the
          environmental monitoring pipeline.
    \item \textbf{Chapter~4} presents experimental results for both
          subsystems, provides full statistical validation of navigation
          performance, reports environmental monitoring outcomes,
          and discusses the limitations of the current implementation.
    \item \textbf{Chapter~5} synthesises the findings into conclusions,
          evaluates the overall contribution of the framework, and
          outlines concrete directions for future research and
          development.
\end{itemize}

Taken together, the chapters present a coherent arc from problem
identification and system design through to rigorous experimental
validation. The reader who follows this sequence will understand not
only what the framework does but also how each design choice was
arrived at, what alternatives were considered, and where the work
can be meaningfully extended. The intention throughout has been
transparency: every hardware decision, every software parameter, and
every experimental result is reported in sufficient detail for the
work to be reproduced independently.